\documentclass[10pt,a4paper]{amsart}
\usepackage[margin=19mm]{geometry}
\usepackage{amsmath,amssymb,mathtools}
\usepackage[scr=rsfs,cal=euler]{mathalfa}
\usepackage{xfrac}
\usepackage[unicode,hidelinks,bookmarks=false]{hyperref}
\newcommand{\ie}{i.e.\ }
\newcommand{\cf}{cf.\ }
\newcommand{\resp}{respectively\ }
\newcommand{\Subsection}{\subsection}
\providecommand{\nobreakdash}{\nobreak\hbox{-}}
\setlength{\emergencystretch}{2em}
\multlinegap0pt
\usepackage{bm}
\usepackage{bbm}
\usepackage{dsfont}
\usepackage{xspace}
\usepackage{cancel}
\usepackage{mathtools}
\usepackage{amsthm}
\makeatletter
\@ifundefined{thm}{\newtheorem{thm}{Thm}}{}
\@ifundefined{prop}{\newtheorem{prop}[thm]{Proposition}}{}
\makeatother
\DeclareMathOperator{\Span}{span}
\title[The geometry of triples of antipodal ideal chambers of affin]{The geometry of triples of antipodal ideal chambers of affine buildings}
\author{Corina Ciobotaru, Corentin Le Bars}
\date{Source: arXiv:2608.04520v1 (CC BY 4.0)}
\begin{document}
\maketitle

\noindent\emph{Source and licence.} The geometry of triples of antipodal ideal chambers of affine buildings. Version arXiv:2608.04520v1 (CC BY 4.0), distributed under the Creative Commons Attribution 4.0 International licence, as recorded in the arXiv licence metadata for this exact version and shown on the abstract page https://arxiv.org/abs/2608.04520v1. Full rights notice: licenses/arxiv-2608.04520-CC-BY-4.0.txt.\par\smallskip

\noindent\textit{Editorial scope.} Counted unit: the Prop whose source label is \texttt{prop:E678} in the article (source0.tex lines 1138--1142, source label \texttt{prop:E678}), delivered together with the article's own complete original proof of it (source0.tex lines 1144--1333, one \texttt{proof} environment of 190 lines). No mathematics is written, repaired or re-derived: statement and proof are the article's own text line for line. Required context retained uncounted: prop:geom\_criterion (lines 654--661). Every definition, display and label of those lines is retained unchanged. Omitted from this package and not redistributed: 1--653, 662--1137, 1143--1143, 1334--1725. Nothing inside a retained excerpt is omitted or altered: every retained body is delivered exactly as the article's lines are written, so no omission receipt of any kind accompanies this package. Citations: the retained excerpts carry no citation command at all, so no bibliography entry of the article is transcribed and no reference is redistributed. Reference adaptation: none was needed and none was made; every reference inside the delivered unit resolves inside it, so no unresolved reference or citation is produced. Printed slips kept literal: no printed word, symbol, hypothesis or number of the counted statement or of its proof was altered. Layout and class: the article's own class is \texttt{amsart} with packages including babel, fontenc, graphicx, amsmath, amsthm, amssymb, mathrsfs, enumitem, xy, tikz-cd, xcolor, hyperref; the wrapper is the lane's pinned \texttt{amsart} editorial wrapper at 10pt on A4, which restates the theorem-like environments the retained text needs and adds the article's own macro definitions that the retained text needs (\texttt{\textbackslash Span}), with extra wrapper packages (bm, bbm, dsfont, xspace, cancel, mathtools). The delivered numbering is the unit's own. Assets: no retained span contains a figure environment, a graphics-inclusion command, a PDF-page inclusion, a table or a TikZ picture; no image file is copied into this package, the package declares no asset dependency and no image file is redistributed with it. Licence: version arXiv:2608.04520v1 of this article is distributed under the Creative Commons Attribution 4.0 International licence, as recorded in the arXiv licence metadata for this exact version and shown on the abstract page https://arxiv.org/abs/2608.04520v1. A full rights notice is delivered at licenses/arxiv-2608.04520-CC-BY-4.0.txt. Characters: every retained excerpt is pure ASCII, so the delivered engine, XeLaTeX with its default Latin Modern text font, reports no missing character.
	\begin{prop}\label{prop:geom_criterion}
		The following conditions are equivalent: 
		\begin{enumerate}
			\item There is no root-flat in the Coxeter complex of the affine Weyl group $(W,S)$ that is perpendicular to a bisector; 
			\item For every root subspace \(U=\Span(T)\subsetneq V\), one has \(\rho\notin U\).
		\end{enumerate}
		
	\end{prop}

	\begin{prop}
		\label{prop:E678}
		The root systems  \(E_6\), \(E_7\) and \(E_8\) violate the conditions of
		Proposition~\ref{prop:geom_criterion}, that is, there are root-flats that are perpendicular to bisectors.
	\end{prop}

	\begin{proof}
		We use Bourbaki's realization of \(E_6\) inside \(\mathbb R^8\), with standard basis \(e_1,\ldots,e_8\).  In this realization the root space is
		\[
		V=\Span_{\mathbb R}\Phi(E_6)
		=
		\{x\in\mathbb R^8 \mid x_6=x_7=-x_8\},
		\]
		so \(\dim V=6\).  The Weyl vector is
		\[
		\rho_{E_6}=(0,1,2,3,4,-4,-4,4).
		\]
		Consider the following five roots of \(E_6\):
		\[
		\begin{aligned}
			\beta_1&=e_1+e_2,\\
			\beta_2&=e_1+e_3,\\
			\beta_3&=e_1+e_4,\\
			\beta_4&=e_2+e_3,\\
			\beta_5&=\frac12(e_1+e_2+e_3+e_4+e_5-e_6-e_7+e_8).
		\end{aligned}
		\]
		The first four roots belong to the \(D_5\)-subsystem \(\{\pm e_i\pm e_j\mid 1\le i<j\le5\}\subset\Phi(E_6)\), and \(\beta_5\) is one of the half-roots in Bourbaki's model.  Let
		\[
		L:=\Span_{\mathbb R}\{\beta_1,\beta_2,\beta_3,\beta_4,\beta_5\}.
		\]
		These five roots are linearly independent.  Indeed, in a relation
		\[
		a_1\beta_1+a_2\beta_2+a_3\beta_3+a_4\beta_4+a_5\beta_5=0,
		\]
		the \(e_6,e_7,e_8\)-coordinates imply \(a_5=0\).  The \(e_4\)-coordinate then gives \(a_3=0\), and the remaining \(e_1,e_2,e_3\)-coordinates give
		\[
		a_1+a_2=0,
		\qquad
		a_1+a_4=0,
		\qquad
		a_2+a_4=0,
		\]
		hence \(a_1=a_2=a_4=0\).  Thus \(\dim L=5\), so \(L\) is a proper root subspace of the six-dimensional root space \(V\).
		
		Finally, a direct calculation gives
		\[
		\rho_{E_6}
		=
		-2\beta_1-\beta_2-\beta_3-\beta_4+8\beta_5.
		\]
		Therefore \(\rho_{E_6}\in L\), and Proposition~\ref{prop:geom_criterion} gives a root-flat perpendicular to a bisector.
		
		
		We turn to $E_8$. In this realization the root system is given by
		\[
		\{\pm e_i\pm e_j\mid 1\le i<j\le 8\}
		\cup
		\left\{
		\frac12\sum_{i=1}^8 \varepsilon_i e_i
		\;\middle|\;
		\varepsilon_i\in\{\pm1\},\;
		\#\{i\mid \varepsilon_i=-1\}\equiv 0 \pmod 2
		\right\}.
		\]
		The simple roots are numbered as in Bourbaki:
		\[
		\begin{aligned}
			\alpha_1
			&=
			\frac12(e_1+e_8-e_2-e_3-e_4-e_5-e_6-e_7),\\
			\alpha_2&=e_1+e_2,\\
			\alpha_3&=e_2-e_1,\\
			\alpha_4&=e_3-e_2,\\
			\alpha_5&=e_4-e_3,\\
			\alpha_6&=e_5-e_4,\\
			\alpha_7&=e_6-e_5,\\
			\alpha_8&=e_7-e_6.
		\end{aligned}
		\]
		Since \(E_8\) is simply laced, the Weyl vector is characterized by
		\[
		\langle \rho_{E_8},\alpha_i\rangle=1,
		\qquad i=1,\ldots,8.
		\]
		Summing the fundamental weights gives
		\[
		\rho_{E_8}=(0,1,2,3,4,5,6,23).
		\]
		
		Set
		\[
		U_8:=\{x\in\mathbb R^8\mid x_1=0\}.
		\]
		The roots
		\[
		\pm e_i\pm e_j,
		\qquad 2\le i<j\le 8,
		\]
		belong to \(\Phi(E_8)\) and form a root subsystem of type \(D_7\). This
		subsystem spans \(U_8\). For instance, the seven roots
		\[
		e_2-e_3,\ e_3-e_4,\ e_4-e_5,\ e_5-e_6,\ e_6-e_7,\ e_7-e_8,\ e_7+e_8
		\]
		are linearly independent and all lie in \(U_8\). Hence \(U_8\) is a proper
		root subspace of the \(E_8\)-root space. Since the first coordinate of
		\(\rho_{E_8}\) is zero, we have
		\[
		\rho_{E_8}\in U_8.
		\]
		Thus \(\rho_{E_8}\) lies in a proper root subspace. By
		Proposition~\ref{prop:geom_criterion}, the affine Coxeter complex of type
		\(\widetilde E_8\) admits root-flats perpendicular to bisectors.
		
		
		We now treat \(E_7\). In Bourbaki's realization, \(E_7\) is the root
		subsystem of \(E_8\) generated by \(\alpha_1,\ldots,\alpha_7\). Its root
		space is
		\[
		V_7
		=
		\{x\in\mathbb R^8\mid x_7+x_8=0\},
		\]
		and
		\[
		\Phi(E_7)=\Phi(E_8)\cap V_7.
		\]
		Again, since \(E_7\) is simply laced, its Weyl vector is characterized by
		\[
		\langle \rho_{E_7},\alpha_i\rangle=1,
		\qquad i=1,\ldots,7.
		\]
		Using the above Bourbaki simple roots, one obtains
		\[
		\rho_{E_7}
		=
		\left(0,1,2,3,4,5,-\frac{17}{2},\frac{17}{2}\right).
		\]
		
		Consider the subspace
		\[
		U_7
		:=
		V_7\cap\{x_1=0\}
		=
		\{x\in\mathbb R^8\mid x_1=0,\ x_7+x_8=0\}.
		\]
		This is a proper subspace of \(V_7\), of dimension \(6\). We claim that it is
		a root subspace. Indeed, the following six roots belong to \(\Phi(E_7)\):
		\[
		\begin{aligned}
			\beta_1&=e_2-e_3,\\
			\beta_2&=e_3-e_4,\\
			\beta_3&=e_4-e_5,\\
			\beta_4&=e_5-e_6,\\
			\beta_5&=e_5+e_6,\\
			\beta_6&=e_7-e_8.
		\end{aligned}
		\]
		They all lie in \(U_7\). Moreover, they are linearly independent. Indeed, a
		linear relation
		\[
		a_1\beta_1+\cdots+a_6\beta_6=0
		\]
		has \(e_7,e_8\)-coordinates \(a_6\) and \(-a_6\), so \(a_6=0\). The remaining
		relation is supported on \(e_2,\ldots,e_6\):
		\[
		\begin{aligned}
			0
			&=
			a_1(e_2-e_3)
			+a_2(e_3-e_4)
			+a_3(e_4-e_5)  \\
			&\qquad
			+a_4(e_5-e_6)
			+a_5(e_5+e_6).
		\end{aligned}
		\]
		Looking successively at the \(e_2,e_3,e_4,e_5,e_6\)-coordinates gives
		\[
		a_1=0,\qquad
		-a_1+a_2=0,\qquad
		-a_2+a_3=0,\qquad
		-a_3+a_4+a_5=0,\qquad
		-a_4+a_5=0.
		\]
		Thus \(a_1=a_2=a_3=a_4=a_5=0\). Hence the six roots
		\(\beta_1,\ldots,\beta_6\) are linearly independent. Since
		\(\dim U_7=6\), they span \(U_7\). Therefore \(U_7\) is a proper root subspace
		of \(V_7\).
		
		Finally, the first coordinate of \(\rho_{E_7}\) is zero and $-\frac{17}{2}+\frac{17}{2}=0$ so $\rho_{E_7}\in U_7$. Thus \(\rho_{E_7}\) lies in a proper root subspace. By
		Proposition~\ref{prop:geom_criterion}, the affine Coxeter complex of type
		\(\widetilde E_7\) admits root-flats perpendicular to bisectors.
		
	\end{proof}

\end{document}
